% !TeX program = xelatex
% Transcription of all nine scans, in source-page order.
% Original medical statements, numbers and English wording are retained.
% Source images: images/pic01.jpg through images/pic09.jpg,
% corresponding to Adobe Scan Oct 21, 2026_1.jpg through _9.jpg.
\ifdefined\settextfont
  \def\TAPStandalone{0}
\else
  \def\TAPStandalone{1}
  \documentclass[11pt,a4paper]{article}
  \usepackage[margin=1.7cm]{geometry}
  \usepackage{amsmath,amssymb,enumitem}
  \usepackage{xepersian}
  \IfFontExistsTF{Amiri}{\settextfont{Amiri}}{\settextfont{FreeSerif}}
  \setlatintextfont{FreeSerif}
  \setlength{\parindent}{0pt}
\fi
\providecommand{\en}[1]{\lr{#1}}
\newcommand{\TAPheader}{{\footnotesize\en{subject 26\qquad Year.\ \_\_\_\quad Month.\ \_\_\_\quad Date.\ \_\_\_}\par}}
\newcommand{\TAPfooter}[1]{\par\medskip{\footnotesize\en{Batus}\hfill #1\par}}
\newenvironment{TAPpoints}{\begin{itemize}[leftmargin=*,itemsep=.2em]}{\end{itemize}}
\newenvironment{TAPenglish}{\begin{latin}\begin{itemize}[leftmargin=*,itemsep=.2em]}{\end{itemize}\end{latin}}
\ifnum\TAPStandalone=1 \begin{document}\fi

% pic01.jpg -- source page 1
\TAPheader
\section*{\en{Trauma to the Abdomen and Pelvis}}
\begin{enumerate}[label=\arabic*.,leftmargin=*,itemsep=.45em]
\item \en{CT is the study of choice in abdominal trauma.}
\item فقط \en{CT} با کنتراست وریدی گرفته می‌شود مگر کنترااندیکاسیون داشته باشیم.
\begin{TAPpoints}
\item پیدا کردن نواحی با کاهش خونرسانی
\item پیدا کردن هماتوم‌ها
\item پیدا کردن خونریزی فعال یا ادرار خارج از محل طبیعی
\end{TAPpoints}
\item اگر قرار است که \en{head CT} هم گرفته شود، باید قبل از تزریق کنتراست انجام شود.
\item کنتراست خوراکی معمولاً جایگاهی ندارد.
\item کنتراست رکتال اگر شک به لسریشن روده داشته باشیم انجام می‌شود.
\item در بدو مراجعه اورژانس یک سونو \en{FAST} برای رد هموپریتوئن انجام می‌شود.
\item شایع‌ترین ارگان‌هایی که در تروماهای بلانت شکم آسیب می‌بینند:
\begin{TAPpoints}\item طحال، کبد، کلیه، مثانه\end{TAPpoints}
\item کبد
\begin{TAPpoints}
\item در مجموع ترومای بلانت و نافذ بیشترین آسیب را می‌بیند.
\item اکثر مرگ‌های ناشی از تروماهای شکم: دلیل آسیب به کبد است.
\item شایع‌ترین محل آسیب: \en{Posterior aspect of the Right Lobe}
\item اکثر آسیب‌های کبد با هموپریتوئن همراه هستند.
\end{TAPpoints}
\end{enumerate}
\TAPfooter{۱}

% pic02.jpg -- source page 2
\clearpage\TAPheader
\begin{enumerate}[label=\arabic*.,leftmargin=*,itemsep=.5em]\setcounter{enumi}{8}
\item ترومای کبد
\begin{TAPpoints}
\item \en{Contrast-enhanced CT Scan}
\item اکثر موارد ترومای کبد: به صورت انتظاری درمان می‌شوند و نیاز به جراحی ندارند.
\end{TAPpoints}
\item حساسیت به ماده حاجب
\begin{TAPpoints}
\item ماده‌ای که امروزه استفاده می‌شود:
\begin{TAPenglish}\item Nonionic low Osmolar Solutions with High Concentration of Iodine\end{TAPenglish}
\item ماده حاجب از طریق کلیه و ادرار دفع می‌شود.
\item در افراد با \en{Cr > 1.5} ماده حاجب موجب \en{ATN} می‌شود.
\begin{TAPpoints}\item دیابت، \en{MM}، بیماران دهیدره\end{TAPpoints}
\item این \en{ATN} عمدتاً برگشت‌پذیر است.
\item در بیمارانی که بیماری زمینه‌ای کلیه دارند ممکن است دائمی شود.
\item واکنش به ماده حاجب وابسته به دوز است.
\item عوارض جانبی: احساس گرما، تهوع و استفراغ، خارش، تحریک موضع تزریق
\begin{TAPpoints}\item عموماً نیاز به درمان ندارند.\end{TAPpoints}
\item بیماران آسمی و یا سابقه حساسیت به تزریق \en{IV} احتمال بالاتری برای واکنش دارند.
\begin{TAPpoints}
\item استروئید، \en{Benadryl}، سایمتیدین
\item قبل، حین یا بعد از تزریق
\end{TAPpoints}
\item حساسیت به صدف ارتباطی ندارد.
\item ریسک آنافیلاکسی \en{0.01--0.04\%}
\end{TAPpoints}
\end{enumerate}
\TAPfooter{۲}

% pic03.jpg -- source page 3
\clearpage\TAPheader
\begin{enumerate}[label=\arabic*.,leftmargin=*,itemsep=.45em]\setcounter{enumi}{10}
\item یافته‌های \en{CT} در ترومای کبد:
\begin{TAPpoints}
\item \textbf{هماتوم ساب کپسولار}
\begin{enumerate}[label=\en{\alph*)},nosep]
\item عموماً پارانشیم کبد در بخش تحتانی را تخت می‌کند.
\item شایع‌ترین محل: انترولترال در لوب راست
\end{enumerate}
\item \textbf{لسریشن}
\begin{enumerate}[label=\en{\alph*)},nosep]
\item شایع‌ترین یافته
\item مارژین نامنظم
\item خطی یا شاخه‌ای
\item معمولاً در پریفرال کبد
\item اگر یک بخش کبد را جدا کند، آن \en{Fracture} می‌گویند.
\end{enumerate}
\item \textbf{هماتوم اینتراهپاتیک}
\begin{TAPenglish}
\item[a)] Focal high attenuation lesions caused by blood
\begin{itemize}[nosep]\item May progress to low attenuation by time\end{itemize}
\item[b)] Mass-like lesions
\end{TAPenglish}
\item \textbf{\en{Wedge-shaped defects}}
\begin{enumerate}[label=\en{\alph*)},nosep]
\item بخش‌های دواسکوله شده پارانشیم کبد که \en{enhance} نمی‌شوند.
\end{enumerate}
\item \textbf{کانتیوژن}
\begin{enumerate}[label=\en{\alph*)},nosep]
\item \en{Minimal parenchymal hemorrhage}
\item مارژین نامشخص
\item تیره‌تر از محیط اطراف
\end{enumerate}
\end{TAPpoints}
\end{enumerate}
\TAPfooter{۳}

% pic04.jpg -- source page 4; continuation of point 11
\clearpage\TAPheader
\begin{TAPpoints}
\item \textbf{پسودوآنوریسم و خونریزی حاد}
\begin{enumerate}[label=\en{\alph*)},nosep]
\item \en{Irregular collection of high attenuation extravasated contrast}
\item معمولاً نیاز به آنژیوگرافی و آمبولیزاسیون دارد یا جراحی
\end{enumerate}
\end{TAPpoints}
\begin{enumerate}[label=\arabic*.,leftmargin=*,itemsep=.6em]\setcounter{enumi}{11}
\item طحال و ترومای آن
\begin{TAPpoints}
\item معمولاً در آسیب‌های دسلریشن دیده می‌شود.
\item اکثر تروماهای طحال درمان محافظه‌کارانه می‌شوند و نیاز به جراحی ندارند.
\item \en{CT Scan} مدالیته انتخابی در بررسی طحال است.
\end{TAPpoints}
\item یافته‌های ترومای طحال در \en{CT}
\begin{TAPpoints}
\item هماتوم ساب کپسولار
\begin{TAPenglish}
\item low attenuation
\item Crescent-shaped collection of fluid
\begin{itemize}[nosep]\item Compresses the normal spleen parenchyma\end{itemize}
\end{TAPenglish}
\item لسریشن
\begin{TAPenglish}
\item Irregular
\item low attenuation
\item Typically transects the spleen
\end{TAPenglish}
\end{TAPpoints}
\end{enumerate}
\TAPfooter{۴}

% pic05.jpg -- source page 5; continuation of point 13
\clearpage\TAPheader
\begin{TAPpoints}
\item \en{Intraparenchymal hematoma}
\begin{TAPenglish}
\item Lacerations filled with blood
\item Are intrasplenic rounded areas of low attenuation
\item May have a mass effect
\begin{itemize}[nosep]\item Enlarge the spleen\end{itemize}
\end{TAPenglish}
\item کانتیوژن
\begin{TAPenglish}
\item Alteration in the normal homogeneous appearance of the spleen.
\item Mottled areas of low attenuation
\end{TAPenglish}
\item مایع یا خون اینتراپریتونئال
\begin{TAPpoints}
\item هموپریتوئن تقریباً در همه آسیب‌ها طحال رخ می‌دهد.
\item می‌تواند به میزان اندک وارد لگن بشود.
\item لزوماً به معنی خونریزی فعال نیست.
\end{TAPpoints}
\end{TAPpoints}
\subsection*{کلیه}
\begin{enumerate}[label=\arabic*.,leftmargin=*,itemsep=.6em]\setcounter{enumi}{13}
\item شایع‌ترین علت ترومای بلانت کلیه
\begin{TAPenglish}\item Motor Vehicle accident\end{TAPenglish}
\item تقریباً همه بیماران با ترومای کلیه: هماچوری
\end{enumerate}
\TAPfooter{۵}

% pic06.jpg -- source page 6
\clearpage\TAPheader
\begin{enumerate}[label=\arabic*.,leftmargin=*,itemsep=.6em]\setcounter{enumi}{15}
\item تصویربرداری انتخابی برای بررسی کلیه
\begin{TAPpoints}
\item \en{Contrast-Enhanced CT Scan}
\item پیلوگرام وریدی و سیستوگرام امروزه خیلی جایگاهی ندارند.
\end{TAPpoints}
\item یافته‌های کلیوی در \en{CT}
\begin{TAPpoints}
\item کانتیوژن
\begin{TAPenglish}
\item Ill-defined
\item patchy low attenuation area in the contrast enhanced CT kidney
\end{TAPenglish}
\item هماتوم ساب کپسولار
\begin{TAPenglish}
\item Crescentic or elliptic densities
\item Compresses the renal parenchyma
\end{TAPenglish}
\item هماتوم پری‌نفریک
\begin{TAPenglish}
\item ill-defined fluid collection surrounding the kidney
\begin{itemize}[nosep]\item Confined by Gerota's fascia\end{itemize}
\end{TAPenglish}
\item لسریشن
\begin{TAPenglish}
\item low attenuation
\item linear or branching
\item May extend to the collecting system or vessels
\end{TAPenglish}
\begin{TAPpoints}\item اگر از کورتکس تا هایلوم: \en{Fracture}\end{TAPpoints}
\end{TAPpoints}
\end{enumerate}
\TAPfooter{۶}

% pic07.jpg -- source page 7; continuation of point 17
\clearpage\TAPheader
\begin{TAPpoints}
\item آسیب‌های عروقی
\begin{TAPpoints}
\item اگر شریانی باشد ممکن است خونرسانی کلیه قطع شود.
\begin{TAPenglish}
\item No contrast enhancement
\item May produce Wedge-shaped defects
\end{TAPenglish}
\end{TAPpoints}
\item آسیب‌های سیستم جمع‌کننده
\begin{TAPpoints}\item خروج ماده کنتراست از لومن\end{TAPpoints}
\end{TAPpoints}
\subsection*{\en{shock Bowel}}
\begin{enumerate}[label=\arabic*.,leftmargin=*,itemsep=.6em]\setcounter{enumi}{17}
\item زمانی که یک ترومای بلانت به شکم
\textit{[عبارت خط‌خورده ناخوانا در اصل]}
\begin{TAPenglish}
\item Severe hypovolemia
\item hypotention
\item Complete reversibility follow resuscitation
\end{TAPenglish}
\item تشخیص در \en{CT}
\begin{latin}
\begin{enumerate}[label=\arabic*),itemsep=.5em]
\item Diffuse thickening of the small bowel wall with increased enhancement
\item Fluid-filled and dilated loops of bowel
\end{enumerate}
\end{latin}
\begin{TAPenglish}
\item Other findings
\begin{itemize}[nosep]
\item Small IVC ($<1$ cm)
% The source writes 6 cm here; retained without clinical correction.
\item Small aorta ($<6$ cm)
\item Decreased perfusion of the spleen
\end{itemize}
\end{TAPenglish}
\end{enumerate}
\TAPfooter{۷}

% pic08.jpg -- source page 8
\clearpage\TAPheader
\subsection*{ترومای لگن}
\begin{enumerate}[label=\arabic*.,leftmargin=*,itemsep=.45em]\setcounter{enumi}{19}
\item ۷۰\% موارد پارگی مثانه ناشی از شکستگی لگن است.
\begin{TAPpoints}\item ۱۰\% از بیماران با شکستگی لگن، پارگی مثانه دارند.\end{TAPpoints}
\item بهترین تصویربرداری برای تشخیص پارگی مثانه
\begin{TAPenglish}\item CT cystogram\end{TAPenglish}
\item پارگی مثانه
\begin{TAPpoints}
\item اکستراپریتونئال
\begin{TAPpoints}
\item معمولاً ناشی از شکستگی لگن
\item ماده حاجب اطراف مثانه می‌ماند.
\end{TAPpoints}
\item اینتراپریتونئال
\begin{TAPpoints}
\item ناشایع‌تر (حدود ۲۰\%)
\item بیشتر مرتبط با ترومای بلانت
\item در کودکان بیشتر دیده می‌شود.
\item \en{Usually at the dome of the bladder}
\item مایع حاجب آزاد حرکت می‌کند.
\end{TAPpoints}
\end{TAPpoints}
\end{enumerate}
\subsection*{آسیب‌های میزراه}
% The word حالب is crossed out and replaced with میزراه in the source heading.
\begin{enumerate}[label=\arabic*.,leftmargin=*,itemsep=.45em]\setcounter{enumi}{22}
\item معمولاً در پی ترومای شدید بلانت لگن که بیشتر در مردان رخ می‌دهد.
\item در \en{Straddle fracture of pelvis} نیز باید در نظر گرفته شود.
\item یافته‌ها
\begin{TAPpoints}
\item هماچوری
\item خون در یورترال مئاتوس
\item ناتوانی در تخلیه ادرار
\end{TAPpoints}
\item تصویربرداری انتخابی:
\begin{TAPenglish}\item Retrograde Urethrography (RUG)\end{TAPenglish}
\begin{TAPpoints}\item باید قبل از نصب کاتتر فولی انجام شود.\end{TAPpoints}
\end{enumerate}
\TAPfooter{۸}

% pic09.jpg -- source page 9
\clearpage\TAPheader
\begin{enumerate}[label=\arabic*.,leftmargin=*,itemsep=.45em]\setcounter{enumi}{26}
\item شایع‌ترین فرم آسیب
\begin{TAPenglish}
\item Rupture of the posterior urethra through the urogenital diaphragm into the proximal bulbous urethra.
\end{TAPenglish}
\end{enumerate}
\subsection*{آسیب‌های شکمی کمتر شایع}
\begin{enumerate}[label=\arabic*.,leftmargin=*,itemsep=.45em]\setcounter{enumi}{27}
\item دیافراگم
\begin{TAPpoints}
\item در حدود ۵\% بیماران ترومایی آسیب می‌بیند.
\item همی‌دیافراگم سمت چپ بیشتر آسیب می‌بیند.
\item بخش پوسترولترال دیافراگم بیشتر آسیب می‌بیند.
\item تشخیص آسیب سمت چپ آسان‌تر است.
\item ممکن است هرنی ارگان‌های شکمی در پی آسیب دیافراگم ایجاد شود.
\item ممکن است بخشی از یک ارگان از هرنی بیرون بزند و گیر بیفتد.
\begin{TAPpoints}\item ممکن است نمای \en{Collar Sign} ایجاد کند.\end{TAPpoints}
\item ممکن است پارگی و یا هرنی به صورت تأخیری رخ دهد.
\end{TAPpoints}
\item پانکراس
\begin{TAPpoints}
\item آسیب پانکراس ناشایع است ولی می‌تواند در تروما نافذ رخ دهد.
\item معمولاً همراه با آسیب به سایر ارگان‌ها است.
\item تصویربرداری انتخابی \en{CT} است.
\item یافته‌ها
\begin{TAPenglish}
\item Peripancreatic fluid
\item Deffuse Enlargement
\end{TAPenglish}
\item واضح‌ترین یافته \en{Fracture} است.
\item از \en{MRCP} برای بررسی مجرا پانکراس استفاده می‌شود.
\item عوارض: سودوسیست، پانکراتیت راجعه، فیستول
\end{TAPpoints}
\end{enumerate}
\TAPfooter{۹}
\ifnum\TAPStandalone=1
  \def\TAPfinish{\end{document}}
\else
  \def\TAPfinish{}
\fi
\TAPfinish
